\chapter[\emj{🧵}~Patrones de Linked List (Reverse, Merge, Reorder)]{\emj{🧵}~Patrones de Linked List}

\begin{dsaquees}

La mayoría de los problemas de listas enlazadas en entrevistas se
reducen a "manipular punteros con cuidado" 🪡. Tres patrones cubren casi
todo:

\begin{itemize}
\item Invertir: darle la vuelta a las flechas, una por una.
\item Fusionar dos listas ordenadas: como cerrar un cierre 🤐, tomando
      cada vez el elemento más pequeño de las dos.
\item Reordenar: partir, invertir e intercalar.
\end{itemize}

El secreto siempre es el mismo: usa un \textbf{dummy node} (nodo falso al
inicio) y muévete con punteros temporales para no perder el resto de la
cadena.

\end{dsaquees}

\begin{dsaotraforma}

Los patrones más rentables sobre listas enlazadas:

\begin{itemize}
\item \textbf{Reverse}: tres punteros (\verb|prev|, \verb|cur|,
      \verb|next|). En cada paso, apunta \verb|cur->next| hacia atrás y
      avanza. $O(n)$ tiempo, $O(1)$ espacio.
\item \textbf{Merge Two Sorted Lists}: un dummy y un puntero \verb|tail|;
      vas enganchando el menor de las dos cabezas. Base del Merge Sort
      sobre listas.
\item \textbf{Reorder / intercalar}: encuentra el medio (fast \& slow),
      invierte la segunda mitad y entrelaza ambas.
\item \textbf{Remove N-th from end}: dos punteros separados por N; cuando
      el de adelante llega al final, el de atrás está justo antes del
      objetivo.
\end{itemize}

El \textbf{dummy node} evita casos borde al modificar la cabeza.

\end{dsaotraforma}

\begin{dsadiagrama}
\begin{verbatim}
REVERSE (3 punteros):
  NULL <- 1    2 -> 3 -> NULL     prev=NULL cur=1
          ^    ^
        prev  cur    guarda next, apunta cur->next=prev, avanza

  Resultado: NULL <- 1 <- 2 <- 3   (head ahora es 3)

MERGE dos listas ordenadas:
  L1: 1 -> 3 -> 5      dummy -> 1 -> 2 -> 3 -> 4 -> 5 -> 6
  L2: 2 -> 4 -> 6      (toma el menor de cada cabeza en cada paso)
\end{verbatim}
\end{dsadiagrama}

\begin{lstlisting}[language=C++]
REVERSE (iterativo, 3 punteros)
- prev=NULL, cur=head. En cada paso: guarda next, cur->next=prev,
  prev=cur, cur=next. Al final prev es la nueva cabeza.

MERGE (con dummy)
- dummy y tail. Mientras ambas listas tengan nodos, engancha el menor
  y avanza esa lista. Al final pega la que sobró.

📝 PSEUDOCÓDIGO (reverse)
──────────────────────────────────────────────────────────────

función reverse(head):
    prev = NULL; cur = head
    MIENTRAS cur != NULL:
        sig = cur.next      // guarda el resto
        cur.next = prev     // voltea la flecha
        prev = cur          // avanza prev
        cur = sig           // avanza cur
    return prev             // nueva cabeza

💻 CÓDIGO C++
──────────────────────────────────────────────────────────────

struct Node { int val; Node* next; };

// Invertir -> O(n) tiempo, O(1) espacio
Node* reverseList(Node* head) {
    Node* prev = nullptr;
    while (head) {
        Node* next = head->next;
        head->next = prev;
        prev = head;
        head = next;
    }
    return prev;
}

// Fusionar dos listas ordenadas
Node* mergeTwoLists(Node* a, Node* b) {
    Node dummy{0, nullptr};
    Node* tail = &dummy;
    while (a && b) {
        if (a->val <= b->val) { tail->next = a; a = a->next; }
        else                  { tail->next = b; b = b->next; }
        tail = tail->next;
    }
    tail->next = a ? a : b;     // engancha lo que sobró
    return dummy.next;
}

// Eliminar el N-ésimo desde el final (two pointers)
Node* removeNthFromEnd(Node* head, int n) {
    Node dummy{0, head};
    Node *fast = &dummy, *slow = &dummy;
    for (int i = 0; i < n; i++) fast = fast->next;   // separa N
    while (fast->next) { fast = fast->next; slow = slow->next; }
    slow->next = slow->next->next;                   // salta el objetivo
    return dummy.next;
}

⏱️ COMPLEJIDAD (Big-O)
──────────────────────────────────────────────────────────────

  Patrón            │ Tiempo  │ Espacio
  ──────────────────┼─────────┼──────────
  Reverse           │ O(n)    │ O(1)
  Merge 2 sorted    │ O(n+m)  │ O(1)
  Reorder           │ O(n)    │ O(1)
  Remove N-th end   │ O(n)    │ O(1)
\end{lstlisting}

\begin{dsaentrevistas}

✅ "Reverse Linked List" (206), "Merge Two Sorted Lists" (21), "Reorder
   List" (143), "Remove Nth Node From End" (19): el póker de listas.
   Domínalos y cubres la mayoría de preguntas de listas.

💡 Dummy node siempre: cuando la operación puede cambiar la cabeza
   (borrar el primer nodo, fusionar), un nodo falso al inicio elimina el
   caso especial y limpia el código.

⚠️ No pierdas la referencia: antes de reasignar \verb|cur->next|,
   SIEMPRE guarda el siguiente en una variable temporal, o pierdes el
   resto de la lista.

🔑 Merge K listas: usa una min-heap con las K cabezas -> O(N log K). Es
   la generalización de "Merge Two Lists" que suele venir después.

\end{dsaentrevistas}

\begin{dsaresumen}

• Reverse: 3 punteros (prev, cur, next); voltea flechas; O(1) espacio.
• Merge ordenadas: dummy + tail, engancha el menor cada vez.
• Reorder: medio (fast/slow) + invertir 2da mitad + intercalar.
• Remove N-th: dos punteros separados por N.
• Usa dummy node cuando la cabeza pueda cambiar.
• Guarda el 'next' antes de reasignar punteros.
════════════════════════════════════════════════════════════════

\end{dsaresumen}
